\chapter*{Kurzfassung}
\addcontentsline{toc}{chapter}{Kurzfassung}

\begin{otherlanguage}{ngerman}
% Keep the complete German summary on one page at the normal text size.
\setstretch{1.45}

Diese Arbeit entwickelt einen Echtzeitregler für einen siebenachsigen
kollaborativen Roboter zur Oberflächenbearbeitung mit einem rechteckigen
Werkzeug auf einer ebenen Fläche. Die kartesische Impedanzregelung erzeugt die
Anpressbewegung und ermöglicht eine passive,
kontaktinduzierte Drehung des Werkzeugs zur Ausrichtung an der Fläche. Dadurch
sollen die Auswirkungen einer Winkelfehlstellung bei Kontaktbeginn reduziert
werden, die andernfalls zu Kantenkontakt und erhöhten Kontaktlasten führen
kann.

Die translatorische und rotatorische Nachgiebigkeit wurde richtungsabhängig zur
Fläche definiert. Ein verschiebbares virtuelles Nachgiebigkeitszentrum koppelt
Anpressbewegung und Rotation. Seine Position relativ zum Werkzeugbezugspunkt
beeinflusst die bevorzugte Ausrichtungsrichtung.

In den Kontaktversuchen wurde die verbleibende Winkelfehlerkomponente um die
erste Flächentangente relativ zur kalibrierten Flächenreferenz untersucht.
Eine Erhöhung der Rotationssteifigkeit
vergrößerte den verbleibenden Fehler um etwa \(276\,\%\).
Die Variation der tangentialen Translationssteifigkeit veränderte ihn um etwa
\(4\,\%\) des Fehlers im Referenzfall.
Die Wirkung der Verschiebung hing von der Richtung des Fehlers bei
Kontaktbeginn ab.
Die kleinsten gemessenen mittleren Fehlerbeträge lagen bei unterschiedlichen
untersuchten Positionen etwa \(53\,\%\) und \(34\,\%\) unter den jeweiligen
Referenzwerten mit dem Nachgiebigkeitszentrum am Werkzeugbezugspunkt.
Die Winkelfehlerkomponente konnte ihr Vorzeichen wechseln. Weitere
kontaktinduzierte Drehung konnte ihren verbleibenden Betrag vergrößern.

In einem getrennten Posehalteversuch wurde die Nullraumregelung untersucht.
Die projizierte Dämpfung verringerte die kumulierte Gelenkbewegung gegenüber
der Einstellung ohne Nullraumdrehmoment um etwa \(25\,\%\).
Die Kombination aus Dämpfung und Singulärwertkonditionierung verringerte
die kumulierte Gelenkbewegung gegenüber der alleinigen Konditionierung
um etwa \(51\,\%\).

\end{otherlanguage}
